\subsection{REPRESENTATION RING}

Let \(\mathfrak{a}\) be a finite dimensional Lie algebra. Let \(\mathcal{F}_{\mathfrak{a}}\) (resp. \(\mathfrak{S}_{\mathfrak{a}}\) ) be the set of classes of finite dimensional (resp. finite dimensional simple) \(\mathfrak{a}\) -modules. Let \(\mathcal{R}(\mathfrak{a})\) be the free abelian group \(\mathbf{Z}^{(\mathfrak{S}_{\mathfrak{a}})}\) . For any finite dimensional simple \(\mathfrak{a}\) -module E, denote its class by {[}E{]}. Let F be a finite dimensional \(\mathfrak{a}\) -module; let \((\mathrm{F}_n,\mathrm{F}_{n - 1},\dots ,\mathrm{F}_0)\) be a Jordan-Hölder series for F; the element \(\sum_{i = 1}^{n}[\mathrm{F}_i / \mathrm{F}_{i - 1}]\) of \(\mathcal{R}(\mathfrak{a})\) depends only on F and not on the choice of Jordan-Hölder series; we denote it by {[}F{]}. If

\[
0 \longrightarrow \mathrm{F} ^ {\prime} \longrightarrow \mathrm{F} \longrightarrow \mathrm{F} ^ {\prime \prime} \longrightarrow 0
\]

is an exact sequence of finite dimensional \(\mathfrak{a}\) -modules, then \([\mathrm{F}] = [\mathrm{F}'] + [\mathrm{F}'']\) .

Let \(\mathrm{F}\) be a finite dimensional semi-simple \(\mathfrak{a}\) -module; for all \(\mathrm{E} \in \mathfrak{S}_{\mathfrak{a}}\) , let \(n_{\mathrm{E}}\) be the length of the isotypical component of \(\mathrm{F}\) of type \(\mathrm{E}\) ; then \([\mathrm{F}] = \sum_{\mathrm{E} \in \mathfrak{S}_{\mathfrak{a}}} n_{\mathrm{E}}. \mathrm{E}\) . If \(\mathrm{F}, \mathrm{F}'\) are finite dimensional semi-simple \(\mathfrak{a}\) -modules, and if \([\mathrm{F}] = [\mathrm{F}']\) , then \(\mathrm{F}\) and \(\mathrm{F}'\) are isomorphic.

Lemma 3. Let G be an abelian group written additively, and \(\varphi : \mathcal{F}_{\mathfrak{a}} \to \mathrm{G}\) a map; by abuse of notation, we denote by \(\varphi(\mathrm{F})\) the image under \(\varphi\) of the class of any finite dimensional \(\mathfrak{a}\) -module F. Assume that, for any exact sequence

\[
0 \longrightarrow \mathrm{F} ^ {\prime} \longrightarrow \mathrm{F} \longrightarrow \mathrm{F} ^ {\prime \prime} \longrightarrow 0
\]

of finite dimensional \(\mathfrak{a}\) -modules, we have \(\varphi(\mathrm{F}) = \varphi(\mathrm{F}') + \varphi(\mathrm{F}'')\) . Then, there exists a unique homomorphism \(\theta: \mathcal{R}(\mathfrak{a}) \to \mathrm{G}\) such that \(\theta([\mathrm{F}]) = \varphi(\mathrm{F})\) for every finite dimensional \(\mathfrak{a}\) -module \(\mathrm{F}\) .

There exists a unique homomorphism \(\theta\) from \(\mathcal{R}(\mathfrak{a})\) to G such that \(\theta([\mathrm{E}]) = \varphi(\mathrm{E})\) for every finite dimensional simple a-module E. Let F be a finite dimensional a-module, and \((\mathrm{F}_{n}, \mathrm{F}_{n-1}, \ldots, \mathrm{F}_{0})\) a Jordan-Hölder series of F; if \(n > 0\), we have, by induction on n,

\[
\theta ([ \mathrm{F} ]) = \sum_ {i = 1} ^ {n} \theta ([ \mathrm{F} _ {i} / \mathrm{F} _ {i - 1} ]) = \sum_ {i = 1} ^ {n} \varphi (\mathrm{F} _ {i} / \mathrm{F} _ {i - 1}) = \varphi (\mathrm{F}).
\]

If \(n = 0\) then \([\mathrm{F}] = 0\) so \(\theta([\mathrm{F}]) = 0\) ; on the other hand, by considering the exact sequence \(0 \longrightarrow 0 \longrightarrow 0 \longrightarrow 0 \longrightarrow 0\) we see that \(\varphi(0) = 0\) .

Example. Take G = Z and \(\varphi(F) = \dim F\) . The corresponding homomorphism from \(\mathcal{R}(\mathfrak{a})\) to Z is denoted by dim. Let c be the class of a trivial a-module of dimension 1, and let \(\psi\) be the homomorphism \(n \mapsto nc\) from Z to \(\mathcal{R}(\mathfrak{a})\) . It is immediate that

\(\dim \circ \psi = \operatorname{Id}_{\mathbf{Z}},\)

so that \(\mathcal{R}(\mathfrak{a})\) is the direct sum of \(\operatorname{Kerdim}\) and \(\mathbf{Zc}\) .

Lemma 4. There exists on the additive group \(\mathcal{R}(\mathfrak{a})\) a unique multiplication distributive over addition such that \([\mathrm{E}][\mathrm{F}] = [\mathrm{E}\otimes \mathrm{F}]\) for all finite dimensional \(\mathfrak{a}\) -modules E, F. In this way \(\mathcal{R}(\mathfrak{a})\) is given the structure of a commutative ring. The class of the trivial \(\mathfrak{a}\) -module of dimension 1 is the unit element of this ring.

The uniqueness is clear. There exists a commutative multiplication on \(\mathcal{R}(\mathfrak{a}) = \mathbf{Z}^{(\mathfrak{S}_{\mathfrak{a}})}\) that is distributive over addition and such that \([E][F] = [E \otimes F]\) for all \(E, F \in S_{a}\) . Let \(E_{1}, E_{2}\) be finite dimensional a-modules, \(l_{1}\) and \(l_{2}\) their lengths; we show that \([E_{1}][E_{2}] = [E_{1} \otimes E_{2}]\) by induction on \(l_{1} + l_{2}\) . This is clear if \(l_{1} + l_{2} \leq 2\) . On the other hand, let \(F_{1}\) be a submodule of \(E_{1}\) distinct from 0 and \(E_{1}\) . Then

\[
[ \mathrm{F} _ {1} ] [ \mathrm{E} _ {2} ] = [ \mathrm{F} _ {1} \otimes \mathrm{E} _ {2} ] \text { and } [ \mathrm{E} _ {1} / \mathrm{F} _ {1} ] [ \mathrm{E} _ {2} ] = [ (\mathrm{E} _ {1} / \mathrm{F} _ {1}) \otimes \mathrm{E} _ {2} ]
\]

by the induction hypothesis. On the other hand, \((\mathrm{E}_1\otimes \mathrm{E}_2) / (\mathrm{F}_1\otimes \mathrm{E}_2)\) is isomorphic to \((\mathrm{E}_1 / \mathrm{F}_1)\otimes \mathrm{E}_2\) . Hence

\[
[ \mathrm{E} _ {1} ] [ \mathrm{E} _ {2} ] = ([ \mathrm{E} _ {1} / \mathrm{F} _ {1} ] + [ \mathrm{F} _ {1} ]). [ \mathrm{E} _ {2} ] = [ (\mathrm{E} _ {1} / \mathrm{F} _ {1}) \otimes \mathrm{E} _ {2} ] + [ \mathrm{F} _ {1} \otimes \mathrm{E} _ {2} ] = [ \mathrm{E} _ {1} \otimes \mathrm{E} _ {2} ],
\]

which proves our assertion. It follows immediately that the multiplication defined above is associative, so \(\mathcal{R}(\mathfrak{a})\) has the structure of a commutative ring. Finally, it is clear that the class of the trivial a-module of dimension 1 is the unit element of this ring.

Lemma 5. There exists a unique involutive automorphism \(\mathrm{X} \mapsto \mathrm{X}^*\) of the ring \(\mathcal{R}(\mathfrak{a})\) such that \([\mathrm{E}]^* = [\mathrm{E}^*]\) for every finite dimensional \(\mathfrak{a}\) -module \(\mathrm{E}\) .

The uniqueness is clear. By Lemma 3, there exists a homomorphism \(X \mapsto X^{*}\) from the additive group \(\mathcal{R}(\mathfrak{a})\) to itself such that \([E]^{*} = [E^{*}]\) for every finite dimensional a-module E. We have \((X^{*})^{*} = X\) , so this homomorphism is involutive. It is an automorphism of the ring \(\mathcal{R}(\mathfrak{a})\) since \((\mathrm{E} \otimes \mathrm{F})^{*}\) is isomorphic to \(E^{*} \otimes F^{*}\) for all finite dimensional a-modules E and F. Q.E.D.

Let \(\mathrm{U}(\mathfrak{a})\) be the enveloping algebra of \(\mathfrak{a}\) , \(\mathrm{U}(\mathfrak{a})^*\) the vector space dual of \(\mathrm{U}(\mathfrak{a})\) . Recall (Chap. II, §1, no. 5) that the coalgebra structure of \(\mathrm{U}(\mathfrak{a})\) defines on \(\mathrm{U}(\mathfrak{a})^*\) a commutative, associative algebra structure with unit element. For any finite dimensional \(\mathfrak{a}\) -module E, the map \(u \mapsto \mathrm{Tr}(u_{\mathrm{E}})\) from \(\mathrm{U}(\mathfrak{a})\) to \(k\) is an element \(\tau_{\mathrm{E}}\) of \(\mathrm{U}(\mathfrak{a})^*\) . If \(0 \longrightarrow \mathrm{E}' \longrightarrow \mathrm{E} \longrightarrow \mathrm{E}'' \longrightarrow 0\) is an exact sequence of finite dimensional \(\mathfrak{a}\) -modules, then \(\tau_{\mathrm{E}} = \tau_{\mathrm{E}'} + \tau_{\mathrm{E}''}\) . Hence, by Lemma 3 there exists a unique homomorphism, which we denote by Tr, from the additive group \(\mathcal{R}(\mathfrak{a})\) to the group \(\mathrm{U}(\mathfrak{a})^*\) such that \(\mathrm{Tr}[\mathrm{E}] = \tau_{\mathrm{E}}\) for every finite dimensional \(\mathfrak{a}\) -module E. If \(k\) denotes the trivial \(\mathfrak{a}\) -module of dimension 1, it is easy to check that \(\mathrm{Tr}[k]\) is the unit element of \(\mathrm{U}(\mathfrak{a})^*\) . Finally, let E and F be finite dimensional \(\mathfrak{a}\) -modules. Let \(u \in \mathrm{U}(\mathfrak{a})\) and let \(c\) be the coproduct of \(\mathrm{U}(\mathfrak{a})\) . By definition of the U-module \(\mathrm{E} \otimes \mathrm{F}\) (Chap. I, §3, no. 2), if \(c(u) = \sum_{i} u_i \otimes u_i'\) ,

\[
u _ {\mathrm{E} \otimes \mathrm{F}} = \sum_ {i} (u _ {i}) _ {\mathrm{E}} \otimes (u _ {i} ^ {\prime}) _ {\mathrm{F}}.
\]

Consequently

\[
\begin{array}{l} \tau_ {\mathrm{E} \otimes \mathrm{F}} (u) = \sum_ {i} \operatorname{Tr} (u _ {i}) _ {\mathrm{E}} \operatorname{Tr} (u _ {i} ^ {\prime}) _ {\mathrm{F}} = \sum_ {i} \tau_ {\mathrm{E}} (u _ {i}) \tau_ {\mathrm{F}} (u _ {i} ^ {\prime}) \\ = (\tau_ {\mathrm{E}} \otimes \tau_ {\mathrm{F}}) (c (u)). \end{array}
\]

This means that \(\tau_{\mathrm{E}}\tau_{\mathrm{F}} = \tau_{\mathrm{E}\otimes \mathrm{F}}\) . Thus, \(\operatorname {Tr}:\mathcal{R}(\mathfrak{a})\to \mathrm{U}(\mathfrak{a})^*\) is a homomorphism of rings.

Let \(a_{1}\) and \(a_{2}\) be Lie algebras, f a homomorphism from \(a_{1}\) to \(a_{2}\) . Every finite dimensional \(a_{2}\) -module E defines by means of f an \(a_{1}\) -module, hence elements of \(\mathcal{R}(\mathfrak{a}_{2})\) and \(\mathcal{R}(\mathfrak{a}_{1})\) that we denote provisionally by \([E]_{2}\) and \([E]_{1}\) .

By Lemma 3, there exists a unique homomorphism, denoted by \(\mathcal{R}(f)\) , from the group \(\mathcal{R}(\mathfrak{a}_2)\) to the group \(\mathcal{R}(\mathfrak{a}_1)\) such that \(\mathcal{R}(f)[\mathrm{E}]_2 = [\mathrm{E}]_1\) for every finite dimensional \(\mathfrak{a}_2\) -module E. Moreover, \(\mathcal{R}(f)\) is a homomorphism of rings. If \(\mathrm{U}(f)\) is the homomorphism from \(\mathrm{U}(\mathfrak{a}_1)\) to \(\mathrm{U}(\mathfrak{a}_2)\) extending \(f\) , the following diagram is commutative

\[
\begin{array}{c c c} \mathcal {R} (\mathfrak {a} _ {2}) & \stackrel {{\mathcal {R} (f)}} {{\longrightarrow}} & \mathcal {R} (\mathfrak {a} _ {1}) \\ \mathrm{Tr} \Big \downarrow & & \Big \downarrow \mathrm{Tr} \\ \mathrm{U} (\mathfrak {a} _ {2}) ^ {*} & \stackrel {{t _ {\mathrm{U} (f)}}} {{\longrightarrow}} & \mathrm{U} (\mathfrak {a} _ {1}) ^ {*}. \end{array}
\]

In what follows we take for \(\mathfrak{a}\) the splittable semi-simple Lie algebra \(\mathfrak{g}\) . The ring \(\mathcal{R}(\mathfrak{g})\) is called the representation ring of \(\mathfrak{g}\) . For all \(\lambda \in \mathrm{P}_{++}\) , we denote by \([\lambda]\) the class of the simple \(\mathfrak{g}\) -module \(\mathrm{E}(\lambda)\) of highest weight \(\lambda\) .

